\begin{abstract}
The base of a supernumerary robotic limb is a person. This project built a
teleoperation system for two seven-degree-of-freedom arms carried on a worn
backpack frame and commanded by a \emph{second} person, gave it a safety
layer appropriate to a manipulator standing beside somebody who is not
commanding it, and ran a first pilot comparing direct control against
shared-autonomy assistance.

A body-worn master supplies position from whichever sensor observes each
component most directly, and classifies all fourteen channels so a failing
sensor degrades the system announcedly rather than silently. Six of fourteen
met the health criterion; failing ones injected \SI{54.5}{\milli\metre} of
unrequested motion. Vertical and fore-and-aft command track correctly;
lateral does not, and no axis mapping fixes it. The two arms share no
reachable point under the as-built mount; a costed mount change raised home
clearance from \SI{0.1610}{\metre} to \SI{0.2202}{\metre}. The task set
solves \num{13927} inverse-kinematics calls with no failures; replacing the
motion generator removed \SI{51.3}{\milli\metre} of unrequested excursion;
and perception on four physical cameras fits a work surface to
\SI{1.43}{\milli\metre} RMS.

Under ethical approval (SETREC 7111986), four operators completed sixteen
sessions comparing direct VR control against shared-autonomy assistance:
assisted sessions averaged shorter completion times and fewer re-grips, with
unchanged hand travel and sub-\SI{5}{\milli\metre} lag. The pilot has no
separate wearer and is pilot-scale; nothing was validated on a loaded
physical arm, and the two-person study this system was built for remains the
target for the work that follows.
\end{abstract}

\renewcommand{\abstractname}{Acknowledgements}
\begin{abstract}
I thank my supervisor, Dr.\ Dandan Zhang, for the direction of this project
and for consistently pushing the work toward what could be measured rather
than what could be claimed. Much of what is useful in this thesis exists
because a result was questioned before it was written down.

I thank my co-supervisor, Dr.\ Etienne Burdet, for the framing that changed
the shape of the project. Treating the wearer and the operator as two people
with different stakes, rather than as one role, is what turned a control
problem into a question worth asking, and it is the reason the workspace and
safety chapters exist in the form they do.

I am grateful to the Human Robotics Group for access to the Dr.\ Octopus
platform, and to everyone who kept two Kinova arms and a backpack frame
available and working.

This project spent a great deal of its time discovering that the instrument
was wrong rather than the system. I am thankful to everyone who responded to a
surprising number with scepticism instead of enthusiasm, which is the only
reason several of them are not in this document as findings.

Finally, I thank my family and friends for their support over the course of
the year.
\end{abstract}
